\documentclass[12pt]{article}
\usepackage{verbatim}
\usepackage{amssymb}
\title{The graph stratification version of the Marcel theorem
prover: documentation and examples}

\author{Lavinia Randall Holmes}

\usepackage{bussproofs}
\usepackage{amssymb}
\usepackage{latexsym}

\begin{document}

\maketitle
\tableofcontents

\section{Introduction}

This document introduces another version of the Marcel theorem
prover for stratified set theory, so-called because it
implements the system of sequent calculus for SF (stratified set
theory with no extensionality axiom but with a set abstraction
operator) defined by Marcel Crabb\'e and shown by him to satisfy
cut elimination (the right rule for equality that we use actually gives the full extensionality of NF, and cut elimination is not known for this system, but this can be corrected if desired).

5/24/2026  documentation is updated to include all user commands available in the software as of that date.

5/26/2026  The embeddded example scripts are temporarily likely to be outdated due to an upgrade in the prover.  The ones that do not work
will be replaced shortly.

This version differs from earlier versions in having no special commands for particular logical operations.  All 
transformations based on rules for particular logical constructions are handled by {\tt left()}, {\tt right()} or {\tt done()} with assistance from
{\tt setunknown} to set particular instances of variables introduced.  It differs from other versions in frankly using full extensionality.  We know that New Foundations itself is consistent.  We could easily install the appropriate modifications of the rules to give SF or NFU.  Unlike the previous version, it does stratification itself automatically (the reasons for this, which are actually the initial trigger of this version of the project, are discussed shortly).  The just previous version used user supplied type indices on bound variables [which we still regard as a good approach].  This version handles defined notions quite differently from earlier versions, and we are rather interested in the results.   This version handles weak stratification issues and variable binding issues generally by keeping track of occurrences of variables invisibly.  This means that in fact only bound variables which are connected to the binding variable of a set abstract need to be typed consistently, which gives a stratification criterion weaker than Marcel's weak stratification (in an inessential way, but it is weaker).   Occurrences
of variables are only identified if they are bound by the same binder.  This has the curious effect that the prover entirely ignores issues
of variable capture: apparent variable capture can happen in the display, but the prover knows which occurrences of variables of the same shape are actually bound in a context with a binder of that shape, and currently keeps track of this without telling the user.  This doesn't seem to actually happen often, but I probably ought to fix the display for user comfort.  The input language is entirely in Polish notation,
which we did for ease of prototyping.  The output language is more familiar-looking.  We probably ought to write a parser for the output language.

We are actively debugging the code:  the example transcripts in the back need to be refreshed now and then and they
will not always be up to date.

This particular version was immediately inspired by
contemplation of the proposal by Ryan Nolan that the
Bellman-Ford algorithm for finding shortest paths in weighted
graphs in which edges might have negative weight could be used
for stratification testing. Our actual proposal owes nothing in
its details to Nolan because the use of that algorithm, while it
works in principle, is a bad idea. We present a rather different
algorithm to be applied to the graph on variables in the
directed graph determined by the atomic formulas in a graph with
weights determined by relative type and order of the variables,
similar to Dijkstra's algorithm. The use of Bellman-Ford is a
bad idea because the execution of this algorithm can be aborted
the second time that any edge is ``relaxed'' (or if an
overestimate of the distance is seen to exist, something that
Bellman-Ford doesnt even notice) with a report of stratification
failure.

It is also motivated by the concurrent and quite independent
study by Thomas Forster and myself of the properties of the
graph on variables in a formula determined by the atomic
formulas in the graph (and in particular the properties of the
notion of connectedness for formulas determined by connectedness
of this graph:  I was already thinking about the graph on variables determined by a formula
and its relation to stratification when I encountered Nolan's suggestion.

This version also follows the path of assigning invisible
occurrence data to every variable in a term or formula.
Occurrences of variables are typed, not the variables
themselves, and this enables direct implemention of weak
stratification (in fact, even weaker than in Marcel's system:
all that is required is that variables connected to the binding
variable in each set abstract have consistent types in relation
to the binding variable.) This has (at least at this initial
stage) the perhaps perverse feature that the definition of
substitution can ignore variable capture: the reason is that the
only instance in which occurrences of variables are regarded as
identical is when they are bound by the same quantifier or set
abstractor. When
substitutions are carried out, the substitution term has all
of its occurrence information displaced to a new range (so that
no identification can exist between any variable
in the substituted term and any binder into whose scope it is
substituted).  Attention to the moral law probably
requires us to install a warning to the viewer when free
variables are introduced into scopes of binders using variables
of the same shape: the distinction is not currently apparent to the viewer
in term or formula display. But the prover knows and does not
make errors due to apparent variable capture.

\section{The language of the prover}

For rapid prototyping, we wrote a parser using Polish notation,
but the display function will look more familiar.

Terms are variables (letters {\tt a-z} with one or more primes
({\tt `} or {\tt *}) affixed) or set abstracts (discussed subsequently). The
primes are prefix in the entry notation (to support the Polish
nature of the parser) and postfix in the display notation.
Variables
have invisible occurrence indices: a new index is generated for
each new variable, but then each occurrence of a variable bound
by a quantifier or set abstractor is coerced to have the same
occurrence index as the binding variable. Occurrence indices are
also manipulated by substitution and by the process of
generating fresh variables in ways to be discussed below.

A set abstract in the input notation takes the form {\tt \{xf},
$\{$ being the abstraction operator, {\tt x} being a variable,
and {\tt f} being a formula. In the display notation it takes
the more familiar form $\{x \mid \phi\}$, $\phi$ being the same
formula but in display notation.

Formulas are of various flavors we now list.

Atomic formulas are of the forms {\tt =tu} (in display form $T =
U$) and {\tt etu} (display form $T \,e \,U$ in the prover: we
might write $T \in U$) where {\tt t} and {\tt u} are terms with
display forms $T$ and $U$ respectively. These are atomic
assertions of equality and membership.

Boolean formulas are of the forms {\tt Cfg} where $C$ is a
connective, one of {\tt \&,V,>,X} in the input language.
Formulas {\tt \&fg, Vfg, >fg, Xfg} take the more familiar forms
$(\phi \, \& \, \psi)$, $(\phi \,V\, \psi)$,
$(\phi\, {\tt ->}\, \psi)$ and $(\phi \,{\tt ==}\, \psi)$
respectively, where {\tt f} and {\tt g} are formulas with
display forms $\phi$ and $\psi$. We may choose to write $(\phi
\wedge \psi)$, $(\phi \vee \psi)$, $(\phi \rightarrow \psi)$ and
$(\phi \leftrightarrow \psi)$ for these familiar logical
operations.

In a separate category due to its arity is the negation {\tt
$\sim$f} which is written $\sim \phi$ in display, where {\tt f}
is a formula with display form $\phi$. We may also write $\neg
\phi$.

Quantified formulas are of the forms {\tt Axf}, {\tt Exf} which
are written $(Ax:\phi)$ and $(Ex:\phi)$ in display (where {\tt
f} is a formula with display forms $\phi$)). We may also write
the usual
$(\forall x:\phi)$ and $(\exists x:\phi)$ here.

The {\tt testt} and {\tt testf} formulas can be used to practice
term and formula entry.

\begin{verbatim}
testt("{x=xx")
'{x | x = x}'
testf("&exyeyz")
'(x e y & y e z)'
\end{verbatim}

The user should be warned that the parser does stratification
checks on set abstracts and will not allow set abstracts which
are not weakly stratified in a suitable sense.
{\small
\begin{verbatim}
testt("{xexx")
x e x
x e x
['stratification failure', [['x', 2], -1, ['x', 2]]]
'!!!'
'!!!'
\end{verbatim}}

Details to be explained later. But notice that the attempt to
present \newline $\{x \mid x \in x\}$ is a failure: one never even sees
it, though one sees its body.

A new feature which has been introduced but which needs debugging is a feature supporting defined terms and formulas.
I did this in a possibly eccentric way which may cause me vast difficulties, or may prove really convenient.  I am not sure.

Capital letters, possibly primed, may be used to represent defined terms and formulas.  The defined terms and formulas
are introduced with {\tt deft} and {\tt deff} commands to be documented.  They are further used by supplying values for variables
in the defined term or formula.  For example, if we define {\tt I} as $\{x:x=a\}$ then we can abbreviate the singleton of $t$ as
$I(a:t)$ ({\tt :atI} in the input language) [$I(a:x)$ actually works correctly:  the reason is that variables derived from definition text are prefixed with at signs ({\tt @}), preventing apparent variable capture].  If we define {\tt C} as $(\forall x:x \in a \rightarrow b \in a)$ then we can represent $t \subseteq u$
as $C(a:t)(b:u)$ (which is {\tt :bu:atC} in the input language;  exchanging order of the arguments would be harmless, as the variable
name controls where the argument is put).   Stratification for these notations is defined (but will no doubt require some debugging):  the term
or formula in a definition must be stratified in the strong sense, so each variable in it has a fixed type, and this information can be used
to build a stratification graph for a term with defined notions in it without unpacking the definitions (the table of relative types of variables in the body of a definition is stored with the definition).

Notice that defined notions may have arbitrary numbers of arguments which may be supplied in any order (since they have associated keys).
The use of defined terms is a bit weird from the standpoint of variable use.  It is best to think of the defined term and the keys for arguments
as bits of text.  One cannot substitute for them and the keys are not variables in the expression.  When a definition with an argument is evaluated,
all unused variables get prefixed with an at sign:  this avoids not actual error but a serious risk of apparent variable capture [a further refinement was added on reflection which averts risk of variable capture in evaluation of let terms, which would have generated error;  it is now secure].  It also warns you when
not all arguments are used.

Merits of this:  there are just two forms of definition, and no more are needed.  Arguments can be supplied in any order, and any number of arguments is supported.  Inhomogeneous operations
are supported (arguments can have arbitrary relative type, though defined terms and formulas must be stratified in the strong sense).   Demerits:  the notation is weird.  The order of the arguments gets reversed between the input and output notation, though this is semantically harmless.

Defined constant terms are permitted (defined terms do not require arguments);  defined constant formulas are not permitted (a defined formula must appear with at least one argument).

In a bit of syntactic sugar, if the first two keys supplied to a defined notation are {\tt a} and {\tt b}, it is displayed in infix form (this works both for terms and for formulas).

\section{Substitutions}

There are two kinds of substitution. One replaces all free
instances of a particular occurrence of a variable in a context
-- this is used for operations on binding contexts (quantifier
rules and comprehension). The other replaces all free variables
with the same surface form (ignoring the occurrence index) in a
context.

In both kinds of substitution, the term substituted in has all
of its occurrence indices displaced above all existing text, so
variables actually cannot be captured when substituted into
binding scopes...but the display
doesn't show this [something we ought to fix]. Later (after
introducing the sequent prover) I'll put in an example of the
weirdness of apparent variable capture.

\section{The stratification algorithm}

The first step of the stratification checker is to build a
directed graph from the formula to be checked. Each atomic
formula generates edges in both directions, with weight equal to
the difference in relative type.
$x=y$ creates edges with weight 0 from $x$ to $y$ and from $y$
to $x$. $x \in y$ creates an edge from $x$ to $y$ with weight 1
and an edge from $y$ to $x$ with weight $-1$. Edges are also
created by
atomic formulas involving set abstracts. For example, $\{x \mid
\phi\} \in z$ creates an edge from $x$ to $z$ of weight 2 and an
edge from $z$ to $x$ of length $-2$, as well as an entire graph
from $\phi$.

Notice that distinct occurrences of free variables are distinct
nodes in the graph. $x \in x$ is a (weakly) stratified formula,
with the two occurrences of $x$ having different relative types.
But occurrences of bound
variables with the same surface form and the same binder are identified: $\{x
\mid x \in x\}$ raises a stratification error.

The stratification algorithm takes as input a node in the
weighted graph and the entire graph. It computes distances from
the selected node to all nodes connected with it. It does this
by starting with a very rough estimate of the distance function
(distance 0 from the selected node to itself [and the trivial
path from the selected node to itself as an extra bit of data]
and infinite distance to all other nodes) and proceeds through
the nodes of the graph improving the estimate at each step. When
a node is encountered, if
there is no finite distance from the selected node already
computed, we move it to the end of the list and proceed to the
next node. If there is a finite distance from the selected node
to the current node, we compute new estimates of the distance
from the selected node to each neighbor of the current node. If
the distance to a neighbor is infinite, we get a finite
estimate, and add that information and a path from the selected
node to the neighbor [both computed in the obvious way] to the
data for that neighboring node. If the distance to a neighbor is
finite, and our new estimate of the distance is the same, we do
nothing to that node. If the distance to a neighbor is finite
and our new estimate is different from the original estimate
(either an overestimate or an underestimate), we can stop the
entire process and report stratification failure, with
supporting data consisting of a cycle of non-zero length constructed from the path to the current node, the
path to the neighboring node previously computed, and the edge
between them. [This is why using the Bellman-Ford algorithm does
not make sense: it will work, but the very first time that a
distance is adjusted from a finite value to a smaller finite
value, all subsequent work is a waste: and the Bellman-Ford
algorithm entirely ignores overestimates, which are also
evidence of failure.] When we arrive at the end of the list of
vertices, we will have computed all finite distances from the
selected node to other nodes, if no stratification error has
been found,
and we can return the partial distance function, which is in
effect a relative type assignment witnessing stratification.
Nodes not connected to the selected node by a path are not
assigned distances.

The system actually uses the stratification algorithm only on
set abstracts [well, it is also used on definition texts, which must be stratified in the strict sense], ensuring that types can be assigned to each bound
variable connected to the binding variable by a path in the
graph in a consistent way. This is even weaker than the usual
``weak stratification'': for example \newline $\{x\mid x=x
\wedge (\exists y:y \not\in y)\}$ is stratified for our
purposes.

\begin{verbatim}

testt("{x&=xxAy~eyy")
'{x | (x = x & (Ay : ~y e y))}'

\end{verbatim}

The function {\tt strattest} can be used to check a formula (not
a term) for stratification.

\begin{verbatim}

testf("e{x=xx{x=xx")
'{x | x = x} e {x | x = x}'

strattest("e{x=xx{x=xx")
[[['x', 27], [1], [['x', 23], 1, ['x', 27]]], [['x', 23], [0],
[['x', 23]]]]

\end{verbatim}

In this formula, there are two different occurrences of $x$, and
the stratification data supplied gives distances for each of
them and a path from the selected node (chosen automatically by
the testing function) to each of the two nodes.

\begin{verbatim}

testf("eu{x&exy&eyzezx")
'u e {x | (x e y & (y e z & z e x))}'
strattest ("eu{x&exy&eyzezx")
[[['z', 52], [-1], [['u', 46], 0, ['x', 47], -1, ['z', 52]]],
[['y', 49], [1], [['u', 46], 0, ['x', 47], 1, ['y', 49]]],
[['x', 47], [0], [['u', 46], 0, ['x', 47]]], [['u', 46], [0],
[['u', 46]]], [['y', 50], [], []], [['z', 51], [], []]]

testf("eu{xEyEz&exy&eyzezx")
(Ey : (Ez : (x e y & (y e z & z e x))))
['stratification failure', [['x', 6], 1, ['y', 7], 1, ['z', 8], 1, ['x', 6]]]
'???'

\end{verbatim}

Here we have two related examples. Neither are stratified in the
strict sense, but the first is weakly stratified, because $y$
and $z$ are free, so different occurrences can be assigned
different types.

In the second formula, where we quantify over $y$ and $z$, there
is no need to run the stratification checker: just trying to
display the formula with the unstratified set abstract in it is a
sufficient disaster to cause the checker to complain.

To assist with reading the output:  a successful stratification is a list of triples, a node followed by the singleton list of the distance to the selected node (or the empty list to signify infinite distance) followed by a path (decorated with weights of edges) from the selected node to the given node.
A report of failure contains two paths from the selected node and an edge between the termini of those two paths, witnessing a change in estimated finite distance.

The emphasis of this work has not been on creating nifty displays for stratification information:  I could improve this, but have not done so so far.

I advise the reader that stratification data is now more crowded.  Installing the definition feature required that I create graphs for terms as well as formulas, and graphs for terms need a base node for the entire term.  So all subterms acquire base nodes, and there are a lot of additional nodes in the graphs (such nodes have names which are numerals).  I clean the base nodes out of stratification display output, but I cannot remove them from graphs [at least, I do not see a way to do it];  they are essential to making the definition feature work.

We note that there is now also a tester for stratified terms rather than formulas ({\tt strattest2})

\section{The sequent prover}

The remainder of the work is development of a user driven sequent prover for NF (which could be adapted to a prover for SF [for which the cut elimination result exists] or NFU by changing the right equality rule).  This parallels my earlier work on versions of the Marcel prover, but it does have interesting features.

The general plan of the sequent prover is that one uses {\tt start s}  to introduce a sequent of the form

$$\begin{array}{c}
\hline \\
P
\end{array}$$

where $P$ is the result of parsing the string $s$.

One then has a limited number of commands as options.

\begin{itemize}

\item The {\tt right()} command applies an appropriate right rule to the first proposition on the right (below the line that is)

\item The {\tt left()} command applies an appropriate left rule to the first proposition on the left (above the line, that is).

\item  The {\tt getleft(n)} command brings the $n$th proposition on the left (above the horizontal bar, actually) to the front.

\item The {\tt getright(n)} command brings the $n$ proposition on the right to the front.

\item The {\tt done()} command invites the prover to recognize that the current sequent is an axiom (that the propositions at the beginnings of the lists on left and right are the same).  Identity of propositions up to $\alpha$-equivalence, ignoring occurrence data, is supported.

The {\tt done()} command is dynamic:  its underlying equality test will make an equation between an unknown variable in one term or formula and the matching term in the other term or formula true by executing the {\tt setunknown} command on the fly.  It does this safely (it tests that the
{\tt setunknown} command will execute without error).  A static equality test will probably be needed eventually, but {\tt done} is currently its only client.

A further, even more dangerous feature has been added.  When the equality function checks whether $x \in u?$ matches $P(x)$
(where $u?$ is an unknown and $P(x)$ is any stratified formula in $x$, it will attempt to set the unknown $u?$ to the value $\{x\mid P(x)\}$.  It does (mod bugs) check that this can actually be done.
This is a kind of higher order matching.  It can be very nice for forestalling the need to manually enter values to be assigned to unknowns.

I have also added reflexivity of equality as an axiom, and the {\tt done()} command handles this.  This may have dynamic effects
if the prover can make the terms in an equation equal by setting unknowns suitably.  The other rules prove reflexivity of equality, but this
should remove lots of junk lines from proofs.

Where the {\tt done} command fails, one might want to use {\tt back()} to restore the previous state, because the prover might have
made undesired global substitutions for unknowns.

\item  The {\tt setunknown(u,t)} command allows one to set the value (globally in the proof) of an ``unknown'' {\tt u} (a free variable introduced by the right rule for the existential quantifier or the left rule for the universal quantifier) with a term {\tt t}.  In a usual presentation, a specific term is supplied when the rule is applied:  here the unity of the {\tt left()} and {\tt right} rules is preserved by separating out the choice of witness.  It may also be useful to delay choosing the witness until it is evident from developments in the proof what might work.   There are limitations on what can replace an ``unknown'' [the term replacing it cannot contain any fresh variable introduced by quantifier rules which is introduced after the unknown which is being replaced].

The command also now allows a value not including any fresh variables to be set for a variable free in the sequent.  This gives theorems and definitions with free variables in them some limited usefulness.

\item  The {\tt Cut(f)} command executes the cut rule with the formula {\tt f}.

\item The {\tt varelim()} command uses an equation at the head of the left list of the sequent to make a global substitution for a fresh variable
when priority conditions on fresh variables permit this.  The equation is then eliminated.  The use of this is to reduce variable clutter, which we recall
is a problem in large proofs which we encountered in older versions of this system.

\item  The {\tt deft(key,term)} command introduces a defined term (defining {\tt key}, a possibly primed capital letter, as {\tt term}) and the {\tt deff(key,formula)} command similarly introduces a defined formula.  These can further be accessed by assigning values to variables: 
in the definiendum:  a defined term can be used as a constant, but a defined formula must have at least one assignment.  I need to provide examples (they are present in the last section).

\item The {\tt savetheorem} command saves the current theorem:  it takes the name to be given to the theorem as its sole argument.

\item  The {\tt theoremcut} command has as its argument the name of a theorem, whose conclusion is introduced as a hypothesis in the current sequent.  This allows saved theorems to be used (examination of the proofs will reveal that this is a hidden use of cut;  one of the sequents in the cut is never seen by the user).

\item The {\tt look()} command lets you look at the current sequent.

\item The {\tt skip()} command lets you move the current sequent to the end of the work queue and go to the next one.  This lets you cycle through all the work to be done and view it (This command is much more sensible than its analogues in older versions, where we used
an actual tree rather than a list of sequents as our data structure).

\item The {\tt back()} command undoes a command, roughly speaking.  A stack of proof states is updated from time to time.  This function
is not entirely reliable [but seems to be improving.]

\item  The {\tt setlog} command sets the name of a log file to which user commands will be appended as they are executed.
This creates scripts which can be used to reconstruct work done under the prover.  I close a log file by setting the log file to a default name ({\tt setlog (``done'')}).  There must be a directory called {\tt logs} in the current directory for a log file to be set up successfully.

\item The {\tt demo()} and {\tt undemo()} commands turn demo mode on and off.  Running a script in demo mode causes a pause after each command (hit any key to continue).

\item The {\tt showdeft}, {\tt showdeff} and {\tt showtheorem} commands display a term or formula definition or theorem named by their string arguments.
The {\tt showtheorems()} command pages through all the theorems (hit any key to continue).  {\tt showdefts()} and {\tt showdeffs()} are analogous.
{\tt showtheproof()} shows the current complete proof all at once, possibly a large object.  {\tt LogTheProof()} does the same thing and also posts it to the log.

\item {\tt makeconstructive()} and {\tt makeclassical()} toggle between classical and intuitionistic logic.  Toggling to intuitionistic logic clears the current proof.  Intuitionistic logic is implemented by causing the {\tt left()} and {\tt right()} commands to discard all right propositions in a sequent except the first one {\em before\/} applying a rule.  Afterwward, it might display more than one proposition on the right:  use {\tt getright()} to decide which one to keep.  Note that the stratified comprehension of this version of Marcel is more liberal than the stratified comprehension used in presentations of intuitionistic NF up to this point.
\end{itemize}

The prover maintains a stack of sequents needing attention.  When the {\tt left()} or {\tt right()} rule is applied, one or two sequents are appended to the end of the proof [which is a list of sequents with pointers (really just lists of line numbers) from sequents to ones justifying them] and one then proceeds to the first sequent on the stack.  Other rules affecting proof state have appropriate effects on the stack.

When every line has a justification, the proof is complete.

The meat of the logic supported is in the functions {\tt leftaction} and {\tt rightaction}. which are used to indicate how to construct the one or two sequents from which a current sequent is derived.

We present the rules.  I have copied these with some editing from a manual for a long-ago version of Marcel.  There is a special additional left rule
for membership in a variable which is described at the end.

This version of Marcel frankly has NF as its logic.  This could be changed straightforwardly.  The logic of this version is classical, but it could very easily be made intuitionistic.

In addition to applying the rules below, the prover also expands defined terms and formulas in a manner which should be documented.

\subsection{Specific Sequent Rules:  Connectives}

We give left rules and right rules (premise and conclusion rules) for the commonly used connectives,
excluding converse implication and xor.  The notations $\Gamma$ and
$\Delta$ represent arbitrary finite sets of propositions.  The
sentence closest to the turnstile on the right or left is the first
sentence in the right or left list in the prover's presentation.  The
premise on the left is the one which is presented first by the prover
when the rule is applied.

The left rule for the biconditional is lazily handled by definitional expansion.

In our experience the rule which it is somewhat difficult to get used
to is the left rule for implication, though with a little thought it
can be seen to express the rule of {\em modus ponens\/}.

\subsection{Axiom}

\begin{center}
\AxiomC{$\Gamma,A\fCenter{\,\vdash\,}A,\Delta$}
\DisplayProof
\end{center}

\subsection{Left rules (Premise rules)}

\begin{center}
\Axiom$\Gamma\fCenter{\,\vdash\,}A,\Delta$
\UnaryInf$\Gamma,\neg A\fCenter{\,\vdash\,}\Delta$
\DisplayProof
\end{center}

\begin{center}
\Axiom$\Gamma,A,B\fCenter{\,\vdash\,}\Delta$
\UnaryInf$\Gamma,A\wedge B\fCenter{\,\vdash\,}\Delta$
\DisplayProof
\end{center}

\begin{center}
\Axiom$\Gamma,A \fCenter{\,\vdash\,}\Delta$
\Axiom$\Gamma,B \fCenter{\,\vdash\,}\Delta$
\BinaryInf$\Gamma,A\vee B\fCenter{\,\vdash\,}\Delta$
\DisplayProof
\end{center}

\begin{center}
\Axiom$\Gamma\fCenter{\,\vdash\,}A,\Delta$
\Axiom$\Gamma,B\fCenter{\,\vdash\,}\Delta$
\BinaryInf$\Gamma,A\rightarrow B\fCenter{\,\vdash\,}\Delta$
\DisplayProof
\end{center}

\begin{center}
\Axiom$\Gamma,A\rightarrow B,B\rightarrow A\fCenter{\,\vdash\,}\Delta$
\UnaryInf$\Gamma,A \leftrightarrow B\fCenter{\,\vdash\,}\Delta$
\DisplayProof
\end{center}

\subsection{Right rules (Conclusion rules)}

\begin{center}
\Axiom$\Gamma,A\fCenter{\,\vdash\,}\Delta$
\UnaryInf$\Gamma\fCenter{\,\vdash\,}\neg A,\Delta$
\DisplayProof
\end{center}

\begin{center}
\Axiom$\Gamma\fCenter{\,\vdash\,}A,B,\Delta$
\UnaryInf$\Gamma\fCenter{\,\vdash\,}A\vee B,\Delta$
\DisplayProof
\end{center}

\begin{center}
\Axiom$\Gamma\fCenter{\,\vdash\,}A,\Delta$
\Axiom$\Gamma\fCenter{\,\vdash\,}B,\Delta$
\BinaryInf$\Gamma\fCenter{\,\vdash\,}A\wedge B,\Delta$
\DisplayProof
\end{center}

\begin{center}
\Axiom$\Gamma,A\fCenter{\,\vdash\,}B,\Delta$
\UnaryInf$\Gamma\fCenter{\,\vdash\,}A\rightarrow B,\Delta$
\DisplayProof
\end{center}

\begin{center}
\Axiom$\Gamma,A\fCenter{\,\vdash\,}B,\Delta$
\Axiom$\Gamma,B\fCenter{\,\vdash\,}A,\Delta$
\BinaryInf$\Gamma\fCenter{\,\vdash\,}A \leftrightarrow B,\Delta$
\DisplayProof
\end{center}
\section{More Sequent Rules}

For conventions on how rules are to be read in general, see the
section on sequent rules for connectives above.  

In the rules which follow, if $\phi$ is a formula, $\phi[t/x]$ is
taken to represent the result of substituting the term $t$ for the
variable $x$ in $\phi$.

\subsection{Rules for Quantifiers}

\subsubsection{Left Rules (Premise Rules)}

\begin{center}
\Axiom$\Gamma,\phi[t/x],(\forall x.\phi)\fCenter{\,\vdash\,}\Delta$
\UnaryInf$\Gamma,(\forall x.\phi)\fCenter{\,\vdash\,}\Delta$
\DisplayProof
\end{center}

\begin{centering}

where $t$ is any term\\
\end{centering}

\begin{center}
\Axiom$\Gamma,\phi[a/x]\fCenter{\,\vdash\,}\Delta$
\UnaryInf$\Gamma,(\exists x.\phi)\fCenter{\,\vdash\,}\Delta$
\DisplayProof
\end{center}

\begin{centering}
where $a$ is a variable not appearing in the conclusion\\
\end{centering}

\subsubsection{Right Rules (Conclusion Rules)}

\begin{center}
\Axiom$\Gamma\fCenter{\,\vdash\,}\phi[a/x],\Delta$
\UnaryInf$\Gamma\fCenter{\,\vdash\,}(\forall x.\phi),\Delta$
\DisplayProof
\end{center}
\begin{centering}
where $a$ is a variable not appearing in the conclusion\\
\end{centering}

\begin{center}
\Axiom$\Gamma\fCenter{\,\vdash\,}\phi[t/x],(\exists x.\phi),\Delta$
\UnaryInf$\Gamma\fCenter{\,\vdash\,}(\exists x.\phi),\Delta$
\DisplayProof
\end{center}
\begin{centering}
where $t$ is any term\\
\end{centering}

\subsubsection{Comments on Quantifier Rules}

In some of the quantifier rules, we have retained the quantified sentence from
the conclusion in the premise.  This is so that we can avoid
formalizing notions of copying and reordering formulas in sequents: a
quantified formula may be reused several times in a proof, and if it
were erased by the application of the rule we would need to copy it
explicitly.  Another advantage is that it preserves precise
equivalence of the conclusion with the conjunction of all the
premises, which is a feature of all the sequent rules of Marcel.

The rules requiring input of a new variable $a$ supply a
computer-generated variable.  In the original version of this prover,
the rules involving a new term $t$ were implemented by separate
commands with the term $t$ as a parameter.  In the current version,
the computer supplies a new ``unknown variable'' which can
subsequently be replaced by a term: the advantage is that the same
command can then handle all basic sequent rules.   The length and readability
of computer generated new variables is greatly improved by having two primes instead of one.

\subsection{Rules for Membership}

\subsubsection{Left Rule}

\begin{center}
\Axiom$\Gamma,\phi[t/x]\fCenter{\,\vdash\,}\Delta$
\UnaryInf$\Gamma,t \in \{x \mid \phi\}\fCenter{\,\vdash\,}\Delta$
\DisplayProof
\end{center}

\begin{centering}
when $\phi$ is stratified (this condition is handled at parse time)\\
\end{centering}

\subsubsection{Right Rule}
\begin{center}
\Axiom$\Gamma\fCenter{\,\vdash\,}\phi[t/x],\Delta$
\UnaryInf$\Gamma\fCenter{\,\vdash\,}t \in \{x \mid \phi\},\Delta$
\DisplayProof
\end{center}

\begin{centering}
when $\phi$ is stratified (this condition is handled at parse time)\\
\end{centering}

\subsection{Rules for Equality}

\subsubsection{Left Rule}

\begin{center}

\Axiom$\Gamma,(\forall v:t\in v \leftrightarrow u\in v)\fCenter{\,\vdash\,}\Delta$
\UnaryInf$\Gamma,t=u\fCenter{\,\vdash\,}\Delta$
\DisplayProof

\end{center}

This is just definitional expansion.  It allows substitution only in stratified contexts, but this is enough for the full strength of substitution as usually presented.

\subsubsection{Right Rule}


\begin{center}
\Axiom$\Gamma\fCenter\fCenter{\,\vdash\,}(Ax.x\in t \leftrightarrow x\in u), \Delta$
\UnaryInf$\Gamma\fCenter{\,\vdash\,}t=u,\Delta$
\DisplayProof
\end{center}

This gives full extensionality.  It can be revised to give SF (use the same Leibniz formulation as in the left rule) or NFU, a bit more elaborate.

To improve the proof theoretic behavior of extensionality, we added a new left rule for membership, tentatively,
which expands $x \in y$ to $$(\forall v:(\forall w:w\in x \leftrightarrow w\in v) \rightarrow v \in y)$$ on the left  (when the other rule does not apply).  We can report that this allows proof of Leibniz equality from coextensionality without cut.  [we revised the exact form of this rule 5/16/2026].

I have also added reflexivity of equality as an axiom.  It is provable from the other axioms, but it is convenient not to have to do it over and over.



\subsection{Global Substitution}

The left rule for the universal quantifier and the right rule for the
existential quantifier require input of a specific term $t$.  What the
prover actually supplies is an fresh variable identified as an ``unknown'', annotated with the list of all fresh variables generated before it.

The variables introduced by the other quantifier rules we call ``arbitrary'' variables.  When unknown and arbitrary variables are considered together, I call them fresh variables.

The {\tt setunknown} command allows the user to set the unknown
variable to a particular term.  There is a restriction on what terms
can replace an unknown variable: the term must have been definable at
the time that the unknown variable was generated, so it
cannot contain any free or unknown variables which were not already declared at the time it was introduced (this data is stored with the unknown variable).

It is worth noting that the display function suffixes arbitrary variables with ! and unknown variables with ?, and it forbids variables of the same shape which are bound from being introduced.  The parser does not require or even understand these suffixes.

An advantage of this approach are that it enables one to delay the
choice of a witness: the later progress of the proof may make it more
evident what witness will work.


\section{Sample proofs}

The ``by lines {\tt [-1]}''  annotations have to do with the fact that the lines are not justified at the time they are displayed in a transcript.  If I displayed the
saved proofs, the list would refer to the numbers of lines justifying the given line,  But the displayed proof file would not show the commands
used to carry out the proof.

We note here (5/6/2026) that we have just installed the ability to log all user commands to an executable file, so  work can be saved and expanded on
(and also run again after revisions of the software.  All example files are now executable Python files generated by the prover as logs of proofs I carried out.  This looks much the same as what I pasted from Python windows in earlier versions, because the logs contain
the state of the display at each step in comments.

\subsection{Inclusion is transitive}

Here is a proof that inclusion is transitive.  The text was generated by the new scripting features of the system.   In this proof,
we used the {\tt LogTheProof()} command to append the final form of the proof to the script file as a comment.

\verbatiminput{docslogs/inclusionistransitivelogfile.py}
\subsection{Equality is transitive}

Here is a proof that equality is transitive, upgraded to use higher order matching.

\verbatiminput{docslogs/neweqtranslogfile.py}

\subsection{Coextensionality implies Leibniz equality}

This example is perhaps a bit strange.  This is a proof that coextensionality implies Leibniz equality:  it uses the weird left rule for membership
statements which I have introduced experimentally, in order to avoid the need to cut.  It is also a nice sequent proof, and examplifies the notation and style of use of the prover, so here it is.


This new version uses the new higher order matching facility.

\verbatiminput{docslogs/coextensionalitylogfile.py}



\subsection{A weird example with unknown variables}

The entry line illustrates the new capability to use spaces, parentheses, brackets, and commas as padding in term and formula input.  There is no balance checking:
parentheses and brackets are just for the person entering the term [to make sure they do it right].

\verbatiminput{docslogs/weirdnesslogfile.py}


\subsection{Atomic inclusion is membership}

This proves the equivalence of $\{x\} \subseteq y$ and $x \in y$.  It involves both term and formula definition.

\verbatiminput{docslogs/atomicinclusionlogfile.py}

\subsection{Projection theorems for the ordered pair}
There may be some wheel spinning in this proof.  It contains definitions of the unordered pair, the ordered pair,
and the first projection operator on ordered pairs, and the proof that the first projection operator works.  This is an enormous text.

\verbatiminput{docslogs/editpairslogfile.py}

\subsection{Foundations of arithmetic}

Also a large file.  Some work on definitions and axioms for arithmetic.

\verbatiminput{docslogs/editarithlogfile.py}




\end{document}